% transcript_zh.tex
% Craig Gentry 在 EUROCRYPT 2021 上关于 FHE 的演讲完整逐字稿（中文翻译版）
% Compile: xelatex transcript_zh.tex

\documentclass[11pt,a4paper]{article}

\usepackage{fontspec}
\usepackage{xeCJK}
\setCJKmainfont[BoldFont=PingFang SC Semibold]{PingFang SC}
\setCJKsansfont[BoldFont=PingFang SC Semibold]{PingFang SC}
\setCJKmonofont{PingFang SC}
\setmainfont{Palatino}
\setsansfont{Helvetica Neue}
\setmonofont{Menlo}

\usepackage[top=2.5cm, bottom=2.5cm, left=2.5cm, right=2.5cm]{geometry}
\usepackage{amsmath,amssymb}
\usepackage{graphicx}
\graphicspath{{images/}}
\usepackage{enumitem}
\usepackage{xcolor}
\usepackage{hyperref}
\usepackage{caption}
\hypersetup{
  colorlinks=true,
  linkcolor=blue!70!black,
  urlcolor=blue!70!black,
  pdfauthor={Whisper Transcription},
  pdftitle={全同态加密十（几）年的发展历程 - 完整逐字稿（中文翻译）},
}

\usepackage{parskip}
\setlength{\parindent}{0pt}
\setlength{\parskip}{6pt plus 2pt minus 1pt}

\usepackage{mdframed}
\newmdenv[
  topline=false, bottomline=false, rightline=false,
  linewidth=3pt, linecolor=blue!40,
  innerleftmargin=12pt, innerrightmargin=12pt,
  innertopmargin=8pt, innerbottommargin=8pt,
  backgroundcolor=blue!3, skipabove=10pt, skipbelow=10pt,
]{myquote}

\usepackage{titlesec}
\titleformat{\section}{\Large\bfseries\sffamily}{}{0em}{}[\vspace{-0.5em}\rule{\textwidth}{0.5pt}]
\titleformat{\subsection}{\large\bfseries\sffamily}{}{0em}{}
\setcounter{secnumdepth}{0}

\setlength{\emergencystretch}{3em}
\tolerance=1000

\newcommand{\keyslide}[2]{%
  \begin{center}
    \includegraphics[width=0.88\textwidth]{#1}
    \captionof{figure}{#2}
  \end{center}
  \vspace{6pt}
}

\newcommand{\timestamp}[1]{%
  \hfill{\scriptsize\color{gray}[#1]}%
}

\begin{document}

\begin{titlepage}
  \centering
  \vspace*{2.5cm}
  {\Huge\bfseries\sffamily 全同态加密十（几）年的发展历程\par}
  \vspace{0.8cm}
  {\LARGE A Decade (or So) of Fully Homomorphic Encryption\par}
  \vspace{1.2cm}
  {\Large\sffamily Craig Gentry\par}
  {\normalsize Algorand Foundation\par}
  \vspace{0.5cm}
  {\normalsize EUROCRYPT 2021 特邀报告\par}
  \vspace{0.3cm}
  {\normalsize 翻译/字幕：刘巍然（学酥）\par}
  \vspace{1.5cm}
  {\large\bfseries 完整逐字稿（中文翻译版）\par}
  \vspace{0.3cm}
  {\small 语音转录：Whisper large-v3-turbo\par}
  \vspace{0.5cm}
  {\small 原始视频：\href{https://www.bilibili.com/video/BV1rY411V7Ko/}{Bilibili BV1rY411V7Ko}\par}
  \vfill
  {\large \today\par}
\end{titlepage}

\newpage
\tableofcontents
\newpage


\section{引言与 FHE 的起源 (00:00 -- 10:00)}

\keyslide{slide_0001.jpg}{幻灯片 00:00}

我们非常高兴也非常荣幸地邀请到 Craig Gentry 作为今天的特邀演讲者。Craig 是 Algorand 基金会的研究员，2009 年在斯坦福大学获得博士学位。众所周知，他撰写了一篇了不起的博士论文，其中构造了第一个全同态加密（FHE）方案。他的论文当然获得了诸多奖项。 \timestamp{00:00}

\keyslide{slide_0002.jpg}{幻灯片 00:39}

他获得了许多奖项，包括 ACM 杰出博士论文奖和 Gödel 奖。Craig 也是 FHE 领域许多重要成果的作者。例如，他是 BGV 方案中的"G"，也是 GSW 方案中的"G"。此外，他在零知识证明、聚合签名和混淆方面也有重要贡献。 \timestamp{00:36}

\keyslide{slide_0003.jpg}{幻灯片 01:19}

今天他的演讲题目是《全同态加密的十（几）年》。请大家一起欢迎 Craig Gentry。非常感谢精彩的介绍和邀请。我将谈谈过去十几年全同态加密的发展历程。当然，这个故事的起点要比这早得多。 \timestamp{01:13}

\keyslide{slide_0004.jpg}{幻灯片 01:58}

故事始于公钥密码学的诞生——Rivest、Adleman 和 Dertouzos 注意到 RSA 方案有一个非常特殊的性质：你可以将密文相乘，这会导致底层明文也相乘。而且你不需要密钥就能做到这一点。因此，实际上你可以在数据保持加密的状态下对其进行计算。这使得数据的处理和数据的访问可以分离。 \timestamp{01:49}

\keyslide{slide_0005.jpg}{幻灯片 02:38}

事实上，早在 1978 年，他们就有了我们今天所说的"加密云计算"的想法。假设有一个用户 Alice，她想使用云计算，但又想保护自己的数据隐私。于是她加密了自己的数据。后来，她对数据有一个查询请求——比如想从云端检索包含某个关键词的文件。 \timestamp{02:26}

她希望的是，尽管数据是加密的，云端仍然能够响应查询、执行查询。为此，加密方案需要有一种"神奇的算法"，能够接收 Alice 发送的函数和她的加密数据，对其进行求值，并产生一个密文，该密文加密的是函数作用于加密内部数据后的结果。这就是同态加密方案所做的事情。 \timestamp{03:02}

\keyslide{slide_0006.jpg}{幻灯片 03:18}

它运行求值函数。此时，云端可以将加密的响应发回给 Alice，她可以解密得到结果。我们对同态加密的要求是：在整个计算过程中，云端在任何中间阶段都不应看到 Alice 数据的任何信息。关于数据不应泄露任何东西——没有部分信息，什么都没有。 \timestamp{03:39}

\keyslide{slide_0007.jpg}{幻灯片 03:57}

关于数据唯一应该泄露的就是数据的长度，因为密文长度必须是数据长度的上界，对吧？所以这一点会泄露，但除此之外，方案中不应泄露任何信息。她的数据不仅要对云端安全，当然也要能抵御任何入侵云端的攻击者或内部攻击等。好的，2009 年，我提出了第一个看似安全的全同态加密方案。抱歉这张照片不太好看。 \timestamp{04:16}

\keyslide{slide_0008.jpg}{幻灯片 04:37}

这是我能找到的唯一一张还不错的照片。基本上，那个方案的核心思想是一种自指的技术，我们称之为"自举"（bootstrapping）。它涉及这样一个问题：如果一个同态加密方案能够用加密的密钥来解密自身会怎样？它接收密钥的加密数据，然后对自身进行解密——就像我第一张幻灯片上的衔尾蛇循环。事实证明，这个想法非常强大。它使我们能够获得全同态加密方案，而且比我 2009 年提出的方案好得多。 \timestamp{04:52}

\keyslide{slide_0009.jpg}{幻灯片 05:16}

这个方案非常慢。你可能想知道它能有多慢。慢到 Butler Lampson 说："我不认为我们有生之年能看到有人使用 Gentry 的方案。"这话让人有点受伤。但回过头来看，他完全正确，因为我们现在拥有的方案比当时好了几百万倍。 \timestamp{05:29}

\keyslide{slide_0010.jpg}{幻灯片 05:56}

许多人做出了贡献。我们现在有四代全同态加密方案，标准化工作也在进行中。这是一段了不起的旅程。以下是演讲其余部分的组织结构。 \timestamp{06:05}

\keyslide{slide_0011.jpg}{幻灯片 06:36}

首先，我想谈谈过去十年的发展——从概念验证到如今可以说是一个可用的工具。我会用一些性能数据和各种应用来"轰炸"大家。虽然我本人偏理论，不是那种会去做 FHE 神经网络评估原型的人。 \timestamp{06:42}

\keyslide{slide_0012.jpg}{幻灯片 07:15}

所以也许我不是最好的推广者，但我觉得有责任更新大家对 FHE 性能的认知。毕竟 2009 年我对 FHE 性能的描述可能过于诚实了，而人们的认知可能就停留在那个时候。所以我觉得我们需要更新认知，看看现在的性能表现。对于"当下"或者说"永恒"的部分，我想描述当前 FHE 方案的工作原理，但我想用一种不同的方式来讲——部分原因是我对典型的"格、LWE、blah blah blah"式的 FHE 演示有点厌倦了。我想把它从格理论中稍微剥离出来，从数学家的角度出发，从同态映射开始，然后描述如何"加固"或"掩盖"同态映射来创建语义安全的加密方案。 \timestamp{07:19}

\keyslide{slide_0013.jpg}{幻灯片 07:55}

最后一部分，我想谈谈未来——是否有希望摆脱我们现有的 FHE 蓝图（使用自举、格和带噪声的密文）？没有什么比摆脱所有这些"包袱"、创造全新的方案更让我高兴的了。那将非常令人兴奋。好的，现在进入第一部分——过去十年 FHE 如何成为一个可用的工具。 \timestamp{07:55}

\keyslide{slide_0014.jpg}{幻灯片 08:35}

如我所说，FHE 经历了四代发展。第一代始于我的初始论文，然后我们有了基于整数的方案和一些简化。2011 年我们实际实现了第一个方案，包括最复杂的自举步骤。但性能确实糟糕透顶。第二代始于 Brakerski 和 Vaikuntanathan 的论文，他们展示了如何基于容错学习（LWE）问题构建 FHE，这是一个成熟的密码学假设。 \timestamp{08:32}

\keyslide{slide_0015.jpg}{幻灯片 09:14}

这非常好。除了安全性上的好处，还有效率上的提升——他们有一种叫做"模数缩减"的技术。在后续论文（通常称为 BGV）中，我们将模数缩减发挥到极致，创建了一个高效得多的方案。参数的增长与你所求值的电路深度相关，而不是度数——深度和度数之间是指数级的差异。所以这是一个很好的性能提升。 \timestamp{09:08}

\keyslide{slide_0016.jpg}{幻灯片 09:54}

第二代的另一个重要部分是 SIMD 操作。基本上，你可以在每个密文中打包大约一千个不同的 \timestamp{09:45}


\section{FHE 的四代发展 (10:00 -- 20:00)}

明文，放在一千个不同的槽位中，并行地对这些槽位进行操作。此外，还有一种方法可以使用自同构来置换不同的槽位，使它们能够相互交互。2012 年，我们有了一个方案，其开销实际上是安全参数的多对数级别，这真的很惊人——因为你会预期任何方案的开销至少是 \timestamp{10:00}

\keyslide{slide_0017.jpg}{幻灯片 10:33}

安全参数级别的，对吧？但它只是安全参数的多对数级别。所以我们当时觉得在性能方面已经做到头了——毕竟怎么可能比多对数更好呢？但事实并非如此。性能在那些最初的年份里确实有了令人难以置信的提升。你可以看到，这里不包括自举——这是一个小小的限定。这些是中等深度的电路，但我们每秒可以执行大约几千个 \timestamp{10:38}

\keyslide{slide_0018.jpg}{幻灯片 11:13}

门操作。你可以看到我们大幅超越了摩尔定律。这是对数刻度。这很好，但不包括自举。如果你想看 2020 年第二代方案的情况，是这样的。这是 HElib 中的自举。如果你看底部高亮的那一行，可以看到每个明文比特的摊销自举时间是 0.9 毫秒。这相当不错。实际上，自举一个密文本身需要 11 秒，你可以在上面一行看到，但每个密文包含大量不同的明文比特。 \timestamp{11:16}

\keyslide{slide_0019.jpg}{幻灯片 11:53}

而且，你不需要在算术电路的每一层都进行自举。那里写着有 17 个可用层级。所以每 17 层自举一次。摊销下来，每个明文比特是 0.9 毫秒。这好多了。你会看到第三代和第四代方案甚至能做得更好。好的，第三代大概始于我与 Sahai 和 Waters 的一篇论文，这是一种简化，也稍微提高了之前方案的效率。 \timestamp{11:55}

\keyslide{slide_0020.jpg}{幻灯片 12:32}

但真正酷的是 Ducas 和 Micciancio 的后续论文，他们展示了可以对该方案进行调整，使其具有非常好的参数。噪声增长不大，基本上可以获得与普通基于格的公钥加密一样好的安全级别。通常 FHE 有非常大的参数，导致安全性大幅损失，但他们证明可以使其与 PKE 一样安全。此外，他们还展示了一种非常高效的自举技术。 \timestamp{12:33}

\keyslide{slide_0021.jpg}{幻灯片 13:12}

在一系列后续论文中，自举变得越来越高效。最后一篇论文中，自举时间不到十分之一秒，那是 2016 年的事。我想此后性能可能还有进一步提升。第三代的自举更快，但这些方案的一个缺点是它们不原生支持密文打包。 \timestamp{13:12}

\keyslide{slide_0022.jpg}{幻灯片 13:51}

在第二代中，你可以在每个密文中塞入大量明文，但第三代并不原生支持这一点。第四代是 CKKS 方案，它操作浮点数，这对神经网络等实际应用非常有用。它极大地提高了效率。以下是 FHE 的各种库。 \timestamp{13:50}

\keyslide{slide_0023.jpg}{幻灯片 14:31}

你可以看到，首先，库非常多。其次，大多数都实现了 CKKS，因为它非常有用，但其他代的方案也有用。事实上，有一种叫做"嵌合 FHE"（Chimeric FHE）的东西，基本思想是在不同的 FHE 方案之间切换非常有用。利用与自举相同的过程——即同态地应用解密函数——你可以从一个 FHE 方案自举到另一个。这样你可以在不同的 FHE 方案之间转换。根据上下文，你可能需要一个在比特上运行良好的第三代方案，或者一个在浮点数上运行良好的第四代方案。 \timestamp{14:29}

\keyslide{slide_0024.jpg}{幻灯片 15:11}

所以你可以在它们之间切换。这张幻灯片来自 Jung-hee Cheon，他做了很多近期的性能改进工作。他展示了自举时间多年来的改进。目前在 CPU 实现上，摊销自举时间为每比特 19 微秒。你可以看到密文相当大。 \timestamp{15:07}

\keyslide{slide_0025.jpg}{幻灯片 15:50}

它们有 262 KB，但包含大量不同的明文，而且你只需要每隔若干层自举一次。所以自举那个非常大的密文只需要 35 秒，摊销下来每比特 19 微秒。他们还做了 GPU 实现，摊销自举时间达到了亚微秒级别，这简直太惊人了。 \timestamp{15:46}

\keyslide{slide_0026.jpg}{幻灯片 16:30}

好的，现在我来谈谈一些应用。再次声明，我是偏理论的人，所以我不了解这些实现的所有细节。我大致筛选了一下，挑出最重要和最令人印象深刻的。FHE 的一个重要应用领域是基因组分析，部分原因是有一个叫 iDASH 的联盟，每年举办一些挑战赛，涉及 MPC 和 FHE，试图在基因组数据上运行隐私保护算法。 \timestamp{16:24}

\keyslide{slide_0027.jpg}{幻灯片 17:10}

密码学家喜欢挑战。每年都有几个团队参加比赛并产出令人印象深刻的成果。这里是 2019 年的全基因组关联研究和基因型填补。他们使用的算法类型包括线性回归、逻辑回归和神经网络。对于 80,000 个 SNP，求值时间为 25 秒。 \timestamp{17:03}

\keyslide{slide_0028.jpg}{幻灯片 17:49}

还有一些训练方面的工作。另一个重要应用领域是神经网络。第一篇真正尝试中等深度神经网络（带有适当非线性激活函数）的论文可能是 2016 年的 CryptoNets 论文。此后，关于神经网络的论文层出不穷。 \timestamp{17:41}

\keyslide{slide_0029.jpg}{幻灯片 18:29}

我选了今年一个我非常喜欢的例子——他们在 100 个不同的人之间进行语音识别，在不到一秒的时间内完成。完全没有简化语音识别算法，就是原样运行，就能在不到一秒内区分 100 个人。另一个应用是隐私信息检索（PIR）。其思想是你想从数据库中检索一个条目，但不告诉服务器你检索的是哪一个。 \timestamp{18:20}

\keyslide{slide_0030.jpg}{幻灯片 19:08}

PIR 的一个必然缺点是服务器必须触及数据库中的每一个条目。因为如果它没有触及某个特定条目，你就会知道你对那个条目不感兴趣。所以服务器在响应 PIR 查询时必须触及数据库中的每一个条目。因此服务器开销是一个大问题，因为它要遍历整个数据库。如果在此基础上的密码学开销还很大，那就很糟糕了。 \timestamp{18:58}

\keyslide{slide_0031.jpg}{幻灯片 19:48}

Sion 和 Carbunar 在 2007 年的论文中说，单服务器 PIR 永远不可能比服务器直接把整个数据库发送给用户更高效。他们的分析基于 Paillier 等数论方案。对于那些方案， \timestamp{19:36}


\section{性能、可用性与标准化 (20:00 -- 30:00)}

也许确实如此，但不久之后，基于格的 PIR 方案进行了重新分析，速度快了几个数量级，表明之前的分析不再适用。对服务器来说，实际响应 PIR 查询比直接传输整个数据库更高效。PIR 方案变得越来越快。我提一下我和 Shai Halevi 的一篇论文，速度非常快，明文和密文的比率是 4/9。所以密文扩展只有大约 2 倍。通信带宽几乎是最优的。 \timestamp{20:00}

\keyslide{slide_0032.jpg}{幻灯片 20:28}

假设查询者需要的信息量有一定下界。我们自己没有做实现，但 Samir Menon 和 David Wu 做了一个初步实现。他们发现，对 100 万个 30KB 的文件进行 PIR 查询需要 86 秒。 \timestamp{20:36}

\keyslide{slide_0033.jpg}{幻灯片 21:07}

你可以将其与未加速的 AES 加密整个数据库进行比较，那大约需要 85 秒。所以服务器响应 PIR 查询的时间基本上与用 AES 加密整个数据库的时间相同，我觉得这真的很惊人。也许你觉得这个比较有点不公平，因为用的是未加速的 AES，但你也可以对 FHE 进行硬件加速，这方面有很大的 \timestamp{21:12}

\keyslide{slide_0034.jpg}{幻灯片 21:47}

投入。DARPA 有一个项目，四年投入 7000 万美元，目标是为 FHE 生产好的硬件加速器。Intel 和 Microsoft 都参与其中，还有其他几个团队。当然，即使没有这个项目，硬件加速也已经在进行中。我之前提到的韩国 Cheon 团队的论文中，摊销时间为每比特 0.423 微秒。 \timestamp{21:48}

\keyslide{slide_0035.jpg}{幻灯片 22:26}

他最近告诉我已经降到了 0.29 微秒。这篇论文实际上是今年发表的，仅在过去几个月内他又提升了约 25\%。几乎可以实时看到性能在提升。好的，FHE 的一个重要因素是可用性。有各种工具包帮助你将未加密数据上的程序转换为加密数据上的程序。 \timestamp{22:24}

\keyslide{slide_0036.jpg}{幻灯片 23:06}

Google 最近开源了一个转译器，将高级代码转换为操作加密数据的代码，引起了关注。目前你需要是密码学专家才能做这些转换，但这个转译器的理念是基本不需要密码学专业知识。在标准化方面，首先，NIST 正在进行后量子密码学的标准化。 \timestamp{23:00}

\keyslide{slide_0037.jpg}{幻灯片 23:46}

他们从 2015 年开始，收到了 82 份提交。后量子密码学是我们希望在大规模量子计算机建成后仍然安全的密码学。在 7 个决赛方案中，有 5 个是基于格的。所以基于格的密码学正在走向标准化，它将长期存在。FHE 将是其中的一部分。我再提一下同态加密标准化研讨会，他们有 \timestamp{23:36}

\keyslide{slide_0038.jpg}{幻灯片 24:25}

不同代同态加密方案的标准草案，还有关于安全级别、API 和应用的持续白皮书。你可以去看看。好的，现在进入演讲的第二部分，我将讲述当前同态加密方案是如何构建的。但我会从一个数学家的角度来讲—— \timestamp{24:13}

首先考虑同态映射，然后考虑如何调整它使其安全。好的，首先我还没有给出任何定义，所以应该定义一些东西。这是公钥加密。你有通常的密钥生成、加密和解密算法。密钥生成和加密是概率算法。我应该在加密算法的输入中加一个小 R 表示随机性，但我没有。请记住 \timestamp{24:49}

\keyslide{slide_0039.jpg}{幻灯片 25:05}

它是概率性的。正确性就是你所想的——如果你加密消息 M 然后解密该密文，对于所有密钥对、消息和加密随机性，你应该得到 M。语义安全是加密方案的正确安全概念，由 Goldwasser 和 Micali 在 1982 年发明。对我们来说，它基本上意味着 0 的加密和 1 的加密应该是不可区分的——对计算受限的攻击者来说应该很难区分。 \timestamp{25:25}

\keyslide{slide_0040.jpg}{幻灯片 25:45}

好的，任何安全的公钥加密方案都必须是概率性的。每个消息必须有多个密文，因为如果每个消息只有一个密文——一个 0 的加密和一个 1 的加密——那你很容易就能破解方案，对吧？因为它是公钥的，你只需加密 0 并与手中的密文比较，就能知道手中的密文加密的是什么。所以任何 PKE 方案都必须是概率性的。 \timestamp{26:01}

\keyslide{slide_0041.jpg}{幻灯片 26:24}

对于同态加密，除了前三个算法外，还有第四个算法，我们称之为 Evaluate（求值）。它应该做的基本上就是我之前展示的加密云计算的那幅图。如果你有一组密文 $c_1$ 到 $c_t$，分别加密 $m_1$ 到 $m_t$，将它们和某个函数 $f$ 一起输入求值函数， \timestamp{26:37}

\keyslide{slide_0042.jpg}{幻灯片 27:04}

输出应该是一个密文，加密的是 $f$ 作用于 $m_1$ 到 $m_t$ 的结果。安全性的定义与基本 PKE 相同，即语义安全。所以同样，每个消息应该有多个密文。用交换图来理解同态的含义可能更容易。你从左上角的一组密文开始。C 代表 \timestamp{27:13}

\keyslide{slide_0043.jpg}{幻灯片 27:43}

密文。你可以沿着向下的箭头解密这些密文得到明文消息，然后对明文消息应用 $f$ 得到另一个消息。或者你可以走"东北通道"——沿着向右的箭头，对密文应用求值函数（内部使用函数 $f$），得到一个密文，解密后会给你那个消息。交换图的主要性质是 \timestamp{27:49}

\keyslide{slide_0044.jpg}{幻灯片 28:23}

无论你走哪条路——东北通道还是西南通道——你都应该得到相同的结果。它是交换的，因为解密函数和函数 $f$ 的应用顺序无关紧要。你可以先应用函数 $f$ 再解密，也可以先解密再应用函数 $f$，应该得到相同的值。 \timestamp{28:26}

\keyslide{slide_0045.jpg}{幻灯片 29:03}

好的，每个同态加密方案都有一个与之关联的函数族，即它实际上能正确工作的函数族。一个函数族可以是算术电路族——基本上是加法和乘法门的分层图。布尔电路可以这样表示，而布尔电路非常强大——基本上一切都可以表示为布尔电路。 \timestamp{29:02}

\keyslide{slide_0046.jpg}{幻灯片 29:42}

所以如果你的同态加密方案能做所有事情、所有函数，我们就称它为全同态加密方案。这就是我们的目标。好的，现在有两种 \timestamp{29:38}


\section{同态加密的两种构造方法 (30:00 -- 40:00)}

看待事物的方式。在本次演讲中，我说有两种开发同态加密方案的方式。第一种是正统的方式——密码学的方式，以 Goldwasser 和 Micali 为代表。你从一个成熟的密码学假设出发，在假设之上构建公钥加密，然后也许你注意到方案已经是同态的，或者你试图通过引入某种加法和乘法操作使其成为同态的。 \timestamp{30:00}

\keyslide{slide_0047.jpg}{幻灯片 30:22}

这种方法的优点是你不太可能得到一个不安全的方案，因为你从一个成熟的密码学假设出发。缺点是你不太可能得到任何同态加密方案，除非你很幸运，你创建的公钥加密方案恰好已经是同态的。我想在这里提出的另一种方式， \timestamp{30:42}

\keyslide{slide_0048.jpg}{幻灯片 31:01}

以 Galois 为代表，是数学的方式。你从一组同态映射出发，将解密解释为属于那个同态映射族，并在此基础上构建同态加密方案。然后你开始考虑安全性——如何"加固"或"修补"或"掩盖"这个同态映射，使其变得安全，从而获得语义安全的加密。如何让我的密文看起来 \timestamp{31:25}

\keyslide{slide_0049.jpg}{幻灯片 31:41}

\keyslide{slide_0050.jpg}{幻灯片 32:21}

与随机不可区分？这种方法的优点是它允许你审视不同的数学结构，考虑它们是否对密码学有用。缺点是，你的方案很可能会被破解。然后你就成了那种提出很多被破解方案的人，不得不承受永恒的耻辱。 \timestamp{32:08}

\keyslide{slide_0051.jpg}{幻灯片 33:00}

我不想说数学的方式就是正确的方式。我认为两种方式都有用，即使在同一个人内部，我认为你也应该时不时考虑两种方法。但我当然不想像 Galois 那样在与密码学家的决斗中丧命。好的，什么是同态映射？基本上，它是两个代数结构之间保持运算的结构保持映射。 \timestamp{32:51}

\keyslide{slide_0052.jpg}{幻灯片 33:40}

用交换图来画这些东西很有用。这是一个同态映射的例子——将 $x$ 映射到 $e^x$。在这种情况下，如果你有 $x$ 和 $y$，从左上角开始，你可以沿着右箭头将它们相加得到 $x+y$，然后向下应用 $\delta$ 得到 $e^{x+y}$。或者从左上角开始，你可以先向下得到 $e^x$ 和 $e^y$，然后将它们相乘得到 $e^{x+y}$。这就是同态映射的一个例子。注意两个结构中的运算不必相同——第一个结构是加法，第二个是乘法。 \timestamp{33:34}

\keyslide{slide_0053.jpg}{幻灯片 34:20}

好的，当你将同态映射与复杂性理论结合时，事情变得有趣了。让我们看一个通用的交换图。假设向下的箭头是同态映射，向右的箭头是运算。现在假设向右的箭头容易计算——它们是可处理的。而向下的箭头——同态映射——很难计算。 \timestamp{34:17}

\keyslide{slide_0054.jpg}{幻灯片 34:59}

除非你有某种陷门、密钥之类的东西。这基本上就是同态加密。它只是同态映射与复杂性理论的一个非常通用的组合。你可以将交换图的顶层视为密文层，你有密文 A、B、C，可以在这个层面执行运算。 \timestamp{35:00}

\keyslide{slide_0055.jpg}{幻灯片 35:39}

但很难沿着向下的箭头到达明文层。所以，如我所说，同态加密基本上就是同态映射加复杂性理论。这个交换图解释了同态映射和解密的关系——解密就是同态映射。 \timestamp{35:42}

\keyslide{slide_0056.jpg}{幻灯片 36:18}

但密钥生成和加密呢？Rothblum 有一个很好的结果：你可以将一个语义安全的私钥同态加密方案（支持模 2 同态加法）通用地转换为语义安全的公钥同态加密方案，能做同样的同态运算。 \timestamp{36:25}

\keyslide{slide_0057.jpg}{幻灯片 36:58}

技术基本上是这样的：你的公钥就是大量 0 的加密和 1 的加密。在加密过程中，你对公钥中密文的一个随机子集进行同态加法，产生一个新的密文来加密你的消息。更具体地说，你有 $k$ 个 0 和 1 的随机加密，分别称为 $c$ 和 $d$。 \timestamp{37:08}

\keyslide{slide_0058.jpg}{幻灯片 37:38}

在加密时，你选择 1 到 $k$ 的一个随机子集 $S$，使得你的消息等于 $|S| \bmod 2$。然后你同态地将对应于 $S$ 中加密 1 的密文相加，对 $S$ 的补集使用加密 0 的密文。你可以看到，如果你把这些密文加起来，加密内容的奇偶性等于 $d$ 的数量，即 $|S|$。 \timestamp{37:51}

\keyslide{slide_0059.jpg}{幻灯片 38:17}

好的，这只是从私钥同态加密产生公钥同态加密的一种通用方法。它很好，因为它也简化了方案的安全性证明——Rothblum 证明了这个加密过程是安全的，你只需要担心证明公钥中密文的不可区分性。所以它简化了可证明安全性。从这个意义上说，同态的方式产生公钥加密方案在精神上与密码学的方式是兼容的。好的，现在我给你一个简单的例子——Goldwasser-Micali 加密方案如何可能由一个用数学方式思考的人来构建。我不知道 Goldwasser-Micali 实际上是怎么做的，但你可以想象作为一个反事实的例子，可能是这样的——你从一个同态映射开始，即 Legendre 符号。 \timestamp{38:34}

\keyslide{slide_0060.jpg}{幻灯片 38:57}

Legendre 符号是从模素数 $p$ 的可逆整数到乘法群 $\{-1, 1\}$ 的同态映射。对此你有这个交换图——解密算法（或向下的箭头）就是应用 Legendre 符号，向右就是乘法。我举个例子：如果你从左上角的 $c$ 和 $d$ 开始（它们是模 $p$ 的数），你可以立即向下应用 Legendre 符号，然后将它们相乘。由于 Legendre 符号的乘法性质，你最终得到的是 $c \times d$ 的 Legendre 符号。你也可以走东北通道——直接将 $c$ 和 $d$ 相乘得到 $c \times d$，然后应用 Legendre 符号。这就是交换图。那么如何从中创建一个安全的加密方案？当然 $p$ 不能公开，因为任何人都能计算 Legendre 符号。 \timestamp{39:17}


\section{经典方案与量子攻击 (40:00 -- 50:00)}

\keyslide{slide_0061.jpg}{幻灯片 39:36}

一个可能的答案——也是 Goldwasser-Micali 的答案——是用一个数 $N$ 来隐藏 $P$，其中 $N = P \times Q$，$P$ 和 $Q$ 未知。然后有二次剩余假设：给定 $N$（而不是 $P$ 和 $Q$），很难区分一个 Jacobi 符号为 1 的模 $N$ 的数 $X$ 是二次剩余还是非二次剩余。这就是用数学方式创建加密方案——我们有一个同态映射，想要加固它，通过用另一个素数 $Q$ 来掩盖它。现在我想用 Rothblum 的加密技术而不是实际的 Goldwasser-Micali 方案来描述加密方案。 \timestamp{40:00}

\keyslide{slide_0062.jpg}{幻灯片 40:16}

\keyslide{slide_0063.jpg}{幻灯片 40:56}

你可以让公钥包含一堆 $X$ 和 $Y$，分别加密 1 或 $-1$。把大量这样的放入公钥。加密过程基本上就是 Rothblum 的方法——你选择 1 到 $K$ 的一个随机子集，取公钥中对应密文的乘积来加密你想要的值。解密就是应用模 $P$ 的 Legendre 符号。 \timestamp{40:49}

\keyslide{slide_0064.jpg}{幻灯片 41:35}

当然，这不是 Goldwasser-Micali 实际的做法。我的观点是他们本可以这样做。他们本可以使用 Rothblum 这种非常通用的技术，利用同态映射来创建语义安全。他们不需要这样做只是因为在这个特定方案中可以做一些简化。但我想强调的是，Rothblum 技术是一种非常通用的方法，可以从同态映射创建语义安全的加密方案。 \timestamp{41:38}

\keyslide{slide_0065.jpg}{幻灯片 42:15}

好的，Goldwasser-Micali 以及 ElGamal 和 Paillier 等其他方案的缺点是：首先，它们不是全同态的。其次，它们不是后量子安全的。量子计算会摧毁这些方案。事实上，量子计算会摧毁所有使用阿贝尔群的同态加密方案，这有点令人沮丧。攻击实际上相当简单。 \timestamp{42:27}

\keyslide{slide_0066.jpg}{幻灯片 42:55}

Watrous 有一个算法，对于阿贝尔群（实际上对于可解群，包括阿贝尔群），他的算法能够计算群的阶。所以如果你使用群同态加密方案，0（或群单位元）的加密形成密文的一个子群。此时你只需要用 Watrous 的攻击来区分加密 0 的子群。如果你想知道某个东西是否是 0 的加密，你只需看它是在这个较小的群中（你可以计算其阶），还是在较大的群中。 \timestamp{43:16}

\keyslide{slide_0067.jpg}{幻灯片 43:34}

\keyslide{slide_0068.jpg}{幻灯片 44:14}

好的，我们看到这些方案不是全同态的。为了获得全同态性，我想转向环同态映射，它允许两种运算——加法和乘法。环同态映射的例子包括模约简（将整数模某个 $n$ 约简）和多项式求值（求值的乘积等于乘积的求值）。 \timestamp{44:05}

\keyslide{slide_0069.jpg}{幻灯片 44:53}

思路是解密算法将是某个秘密的同态映射，它与加法和乘法都交换。但这种方法立即就有一个问题——量子再次来袭。如果你看 0 的加密，如果解密真的是这个环同态映射，那么 0 的加密在密文环中形成一个理想。它是同态映射的核，特别是它是密文群的一个加法子群，Watrous 的求阶攻击再次适用。所以这不可能是后量子安全的。 \timestamp{44:55}

\keyslide{slide_0070.jpg}{幻灯片 45:33}

更糟的是，基本上几乎总是存在一个平凡的线性代数攻击——你只需将密文集合视为一个向量空间，0 的加密将形成一个子空间，用线性代数很容易区分两者。更糟的是——也许有某种方法可以修补这种方法，这就是我们接下来要尝试的。1978 年，Rivest、Adleman 和 Dertouzos 有一个提案，基本上是一种环同态方法，使用合数模数，秘密是一个因子 $p$。加密就是模 $p$ 约简。问题是如何隐藏 $p$，他们用 $N$ 来做，但在这个设定中这不足以保证语义安全，因为 0 的加密都是 $p$ 的倍数。 \timestamp{45:44}

\keyslide{slide_0071.jpg}{幻灯片 46:13}

在他们的方案中，消息 $m$ 基本上被加密为一个模 $p$ 等于 $m$ 的数。所以 0 的加密是 $p$ 的实际倍数，只需取它们的 GCD 就能恢复密钥。他们当时就知道选择明文攻击，但仍然好心地提出了这些方案，我们都应该感谢他们，因为它实际上与我接下来要讨论的、今天使用的方案非常相似。好的，修补这个方案的一个想法是：也许它能给出一个单向加密方案，如果明文有足够的熵，也许它能工作。这基本上就是我们要做的。 \timestamp{46:33}

\keyslide{slide_0072.jpg}{幻灯片 46:52}

\keyslide{slide_0073.jpg}{幻灯片 47:32}

我们将通过给消息添加噪声或熵来隐藏秘密 $p$，以击败 GCD 攻击。最终你得到的是近似 GCD 假设：如果你采样大量数 $x$，每个都非常接近 $p$ 的倍数但不是精确的倍数，那么这些整数（假设是）与相同大小的随机整数不可区分。这就是近似 GCD 假设。我们可以用它来创建同态加密方案。 \timestamp{47:22}

\keyslide{slide_0074.jpg}{幻灯片 48:11}

方案如下：密钥是一个大整数 $p$，然后我们按照 Rothblum 的方式创建公钥，从一个随机的近似 GCD 实例开始——这些整数 $x$ 都非常接近 $p$ 的倍数。为简单起见，我说它们模 $p$ 的余数是偶数。然后我们有 $y$ 和 $z$，其中 $y$ 就是初始的 $x$（模 $p$ 余数为偶数），$z$ 是 $x$ 加 1（模 $p$ 余数为奇数）。 \timestamp{48:11}

\keyslide{slide_0075.jpg}{幻灯片 48:51}

加密时，你使用 Rothblum 的方法——选择这些密文的一个随机子集，使得你包含的 $z$ 的数量加起来：如果要加密 1 则为奇数，如果要加密 0 则为偶数。解密时，首先模 $p$ 约简，应该得到这个值，然后模 2 约简。你应该得到 $|S|$，即被加密的消息。对于同态运算，你只需将这些整数密文相加和相乘。 \timestamp{49:00}

\keyslide{slide_0076.jpg}{幻灯片 49:31}

这隐式地将模 $p$ 的余数相加和相乘，进而隐式地将模 2 的值相加和相乘。 \timestamp{49:50}


\section{噪声管理与 LWE (50:00 -- 60:00)}

\keyslide{slide_0077.jpg}{幻灯片 50:10}

安全性基于 Rothblum 的过程和近似 GCD 假设。这里有两个问题：噪声在增长。这个等式 $C \bmod P$ 等于 $2R$ 那些东西——如果 $R$ 大于 $P$ 并且绕回模 $P$，这个等式可能不成立。那时你就会有解密错误。第二个问题是，当你将这些整数相乘时，会得到超大的整数。 \timestamp{50:00}

密文的比特数在爆炸式增长。好的，我跳过这张幻灯片。我想给你另一个使用多项式求值的环同态方法的例子。这里我们令 $R$ 为模 $Q$ 的整数环。同态映射 $\delta$ 就是在某个秘密点 $S$ 处求值多项式。 \timestamp{50:12}

好的，我们从这个同态映射出发，想要创建一个同态加密方案。如何使其安全？显然我们必须隐藏 $S$。但这并不真正有效，因为 0 的加密都是以 $S$ 为公共根的多项式。 \timestamp{50:24}

\keyslide{slide_0078.jpg}{幻灯片 50:50}

人们常说计算 Gröbner 基是困难的。但在这种情况下，如果你不想让密文中的单项式数量增长，如果你的密文中有可管理数量的单项式，那么你可以在这些单项式上产生一个向量空间。你会有一个密文多项式的向量空间。 \timestamp{50:36}

0 的加密将在该空间中形成一个子空间。线性代数可以轻松区分 0 的加密和其他东西的加密。Fellows 和 Koblitz 在 1993 年基本上就有这个方案作为基本加密方案。他们没有真正建议同态运算，因为我提到的问题——当你加和乘密文时， \timestamp{50:48}

单项式会爆炸式增长，变得低效。所以我们想修复这个方案。怎么做？我们要添加噪声和熵来击败线性代数。 \timestamp{51:00}

这里我们自然地到达了容错学习（LWE）假设。用多项式的语言来说，它是这样的：在秘密 $s$ 处求值为非常小的值（模 $q$）的线性多项式 $f_i(x)$， \timestamp{51:12}

\keyslide{slide_0079.jpg}{幻灯片 51:30}

与可能求值为大值的随机线性多项式不可区分。 \timestamp{51:24}

使用这个假设，我们产生了 Brakerski-Vaikuntanathan 同态加密方案。密钥只是一个随机点。公钥可以按照 Rothblum 的方式产生—— \timestamp{51:36}

基本上产生 $g$ 和 $h$ 作为求值为小数的多项式，对于 0 的加密为偶数，对于 1 的加密为奇数。然后我们对公钥中密文的一个随机子集进行同态加法， \timestamp{51:48}

\keyslide{slide_0080.jpg}{幻灯片 52:09}

产生一个新的密文，其噪声的奇偶性等于你想加密的消息。 \timestamp{52:00}

解密就是——首先在 $s$ 处求值密文多项式，模 $q$。 \timestamp{52:12}

应该等于关于 $e$ 的那个表达式。如果是这样，你可以模 2 约简， \timestamp{52:24}

\keyslide{slide_0081.jpg}{幻灯片 52:49}

恢复 $|S|$，即 $m$。好的，同态映射就是密文多项式的加法和乘法， \timestamp{52:36}

它对内部的消息做了正确的事情。好的，和整数方案一样， \timestamp{52:48}

我们有类似的问题。首先，噪声在增长。其次， \timestamp{53:00}

当你乘密文多项式时，度数会上升。它们不再是线性多项式了。那么如何解决这个问题？首先， \timestamp{53:12}

\keyslide{slide_0082.jpg}{幻灯片 53:28}

这是如何重新线性化的。我们有两个线性密文多项式，想要将它们相乘。 \timestamp{53:24}

乘完后得到一个二次多项式。希望如果你在 $s$ 处求值这个二次多项式， \timestamp{53:36}

你会得到正确的结果——$e_1 \times e_2 \bmod q$，模 2 约简后得到 $m_1 \times m_2$。但密文多项式是二次的。 \timestamp{53:48}

\keyslide{slide_0083.jpg}{幻灯片 54:08}

我们想消除二次项。如何减去它们？我们用"略微二次"的多项式来扩充公钥。 \timestamp{54:00}

这些多项式大部分是线性的，只有一个二次项。例如， \timestamp{54:12}

这个多项式有一个二次项 $2^t \cdot x_j \cdot x_k$。你可以用这些略微二次的多项式来减去 \timestamp{54:24}

\keyslide{slide_0084.jpg}{幻灯片 54:48}

密文多项式中的所有二次项。你把所有二次项都减掉。 \timestamp{54:36}

最终你得到一个再次线性的密文，并且加密了正确的东西。我没有特别提到， \timestamp{54:48}

但这里的误差都是偶数，所以当你减去 \timestamp{55:00}

\keyslide{slide_0085.jpg}{幻灯片 55:27}

这些多项式时，你保持了误差模 2 的奇偶性，使其加密正确的东西。 \timestamp{55:12}

好的，第二个问题是自举。你有一个密文多项式， \timestamp{55:24}

其噪声越来越大。如果你再对它执行更多操作，它会绕回模 $Q$，你会得到解密错误。 \timestamp{55:36}

如何将误差减小到可管理的大小？一个想法是尝试通过减去 \timestamp{55:48}

\keyslide{slide_0086.jpg}{幻灯片 56:07}

精心选择的 0 的加密来降低噪声。你不改变加密的内容，但这并不真正有效，因为你不知道如何减去正确的噪声。 \timestamp{56:00}

你不知道密文内部的噪声。所以我们还没有让这种方法奏效。 \timestamp{56:12}

第二个想法是——你有解密算法，解密和密钥本来就是用来去除噪声的。 \timestamp{56:24}

\keyslide{slide_0087.jpg}{幻灯片 56:46}

这就是它的全部意义。有没有某种方法可以利用它们来降低加密方案内部的噪声水平？ \timestamp{56:36}

这是它的工作原理。这是我们加密方案的交换图。 \timestamp{56:48}

事实上，这个交换图对某个函数族 $F$ 有效。它还不是所有算术电路。 \timestamp{57:00}

\keyslide{slide_0088.jpg}{幻灯片 57:26}

它是"有限同态"的——对一定度数以内的多项式有效，然后噪声变得太大 \timestamp{57:12}

就不行了。但它对某个函数族 $F$ 有效。我们的目标是产生某个新鲜的密文， \timestamp{57:24}

加密相同的消息。我们从某个加密消息 $M$ 的密文 $C$ 开始，想要"刷新" \timestamp{57:36}

密文使其噪声更小。我们想使用我们的交换图。 \timestamp{57:48}

\keyslide{slide_0089.jpg}{幻灯片 58:06}

这就是交换图真正大放异彩的地方。我们想完成交换图。怎么做？ \timestamp{58:00}

左下角放什么？我声称，由于消息 $M$ 是隐藏的，而这个刷新密文的自举过程 \timestamp{58:11}

是公开的，执行刷新的人对 $M$ 一无所知， \timestamp{58:24}

\keyslide{slide_0090.jpg}{幻灯片 58:45}

除了 $C$ 是加密 $M$ 的密文。基本上，你能放在这里来完成交换图的唯一东西 \timestamp{58:36}

就是 $C$，因为 $C$ 是我关于 $M$ 唯一知道的东西。一旦我把 $C$ 放在那里， \timestamp{58:48}

这意味着函数 $F$ 必须是解密函数，因为那是我知道的从 $C$ 到 $M$ 的唯一方法。 \timestamp{59:00}

\keyslide{slide_0091.jpg}{幻灯片 59:25}

一旦我确定了 $C$ 在这里，上面左上角的密文—— \timestamp{59:12}

这是解密箭头，对吧？所以这些必须是密文各比特的加密。 \timestamp{59:24}

一旦我们对这些密文应用求值函数， \timestamp{59:36}

使用我们确定为解密函数的函数 $F$， \timestamp{59:48}


\section{FHE 的未来 (60:00 -- 65:22)}

\keyslide{slide_0092.jpg}{幻灯片 60:05}

我们在这里得到的是密文 $C^*$，它是求值函数应用于内部的解密函数和密文 $X_i$ 的结果。我不会完全按这种方式计算，我会稍作修改——解密函数中我们不知道密钥，对吧？所以我们不能精确求值这个函数。相反，我们可以将其视为 \timestamp{60:00}

\keyslide{slide_0093.jpg}{幻灯片 60:44}

解密函数，我们知道如何用加密的密钥和加密的密文比特来求值。这样我们基本上得到同样的东西——密文 $C^*$，它加密 $M$——前提是这个解密函数在方案能够求值的函数族 $F$ 中。 \timestamp{60:50}

\keyslide{slide_0094.jpg}{幻灯片 61:24}

所以这个交换图中，一切都是预先确定的。一旦你决定要通过利用交换图来刷新密文——这本质上是同态加密方案的精髓——你想使用交换图并完成它，基本上一切都是确定的。自举似乎在某种意义上是唯一真正有效的 \timestamp{61:41}

\keyslide{slide_0095.jpg}{幻灯片 62:03}

\keyslide{slide_0096.jpg}{幻灯片 62:43}

刷新密文的过程。好的，我想我已经超时了，但让我看看。我想说的是，我们最终实现了这个交换图。如果你将密文空间视为一个代数结构，它非常奇怪。 \timestamp{62:32}

\keyslide{slide_0097.jpg}{幻灯片 63:23}

它并不保留明文空间的原始环结构。它上面有二元运算，这些运算是交换的，但不是结合的或其他什么。所以它具有所谓"原群"（magma）的结构——在代数结构中，除了交换性之外，它是最大程度上无结构的。 \timestamp{63:23}

\keyslide{slide_0098.jpg}{幻灯片 64:02}

人们不禁想知道，这是否是全同态加密方案的固有特征——密文空间和通过求值算法执行的运算必须在这个意义上非常无结构。好的，我没有真正留出时间来讨论要探索的新同态映射。我试着在两分钟内快速过一下这些幻灯片。 \timestamp{64:14}

\keyslide{slide_0099.jpg}{幻灯片 64:42}

如我所说，人们不禁想知道密文是否固有地需要极度无结构才能获得全同态加密方案。首先，当密文空间是可解群时，我们看到了一个很大的问题——存在量子攻击。非可解群比较有趣，因为你可能利用 Barrington 定理获得一种同态映射——Barrington 证明了如何使用非阿贝尔群 \timestamp{65:05}

来求值电路，你可以做类似的事情。所以非可解群在这个意义上很有趣，但使用它们有各种问题。量子子群攻击可能仍然适用——我们可能只需取这些非可解群的一个阿贝尔子群，仍然能够区分大群和子群。还有各种其他问题。 \timestamp{65:55}

另一个问题是这些群通常是线性的。线性群基本上意味着矩阵群——它同构于一个矩阵群。表示论基本上就是对矩阵群的研究，它非常强大。 \timestamp{66:46}

例如 Braid 群——首先它们在很多小方面被破解了，但最终被彻底破解是因为利用表示论发现它们是线性的。是的，有各种各样的问题。我仍然对基于多变量密码学的新 FHE 方案抱有希望——你的密文空间是某种原群，也许你通过做某种奇怪的逆运算来执行操作，应用 $F$ 和 $F^{-1}$——这在多变量系统中很常见。 \timestamp{67:37}

这样你会在密文空间中以伪随机的方式跳来跳去，有点像基于格的方案所做的。但我还没有完全探索这个方向。 \timestamp{68:28}

不过，至少为了缩写的缘故，有一个 HFE-FHE 方案也会很酷。总结一下，已经取得了巨大进展，希望未来会有更多突破。谢谢大家。 \timestamp{69:19}


\end{document}
